\documentclass[11pt]{article}

\usepackage[margin=1in]{geometry}
\usepackage{amsmath,amssymb}
\usepackage{booktabs}
\usepackage{graphicx}
\usepackage{hyperref}
\usepackage{xcolor}
\usepackage{enumitem}

\title{\vspace{-1em}Zero-Retraining Domain-Size Transfer for Spatiotemporal
Chaos, and a Damping-Controlled Route to It in Two Dimensions}
\author{Dalton Jones}
\date{\today}

\newcommand{\DKY}{D_{\mathrm{KY}}}

\setlength{\parskip}{0.3em}
\setlength{\parindent}{0pt}
\raggedbottom

\begin{document}
\maketitle
\vspace{-2em}

\noindent\textit{Summary memo. Full architecture, methodology, every
number below, and the complete experiment log are in the companion
document, \texttt{ks\_local\_field\_ltransfer.tex}, sent alongside this
one; section references below (\S) point there.}

\section{Architecture, in one paragraph}

Instead of encoding a PDE's state into one global latent vector, we
encode it into a \emph{field} of local latent sites laid out on a
periodic lattice that mirrors the physical domain: a stack of circular
convolutions (translation-equivariant by construction, with a fixed,
non-learned gauge-anchor channel per site) maps the field to
$P$ sites $\times$ $c$ channels, and a weight-shared local propagator
advances every site using only its neighbors. Because every learned
weight is a per-site operator, nothing in the model is tied to the
domain size $L$ except the physical patch length each site represents.
Growing $L$ at fixed resolution therefore only adds \emph{more sites},
never a differently-shaped weight matrix --- so a checkpoint trained at
one $L$ can be loaded, unchanged, into a model with more sites and run
at a larger $L$ with zero retraining. (\S2, \S\,3.2)

\section{Result 1: clean zero-shot transfer in 1D Kuramoto--Sivashinsky}

Trained once at $L{=}100$ ($\DKY{=}22.14$), the model was evaluated
zero-shot at $2\times$ and $4\times$ the domain, against freshly
computed ground-truth Lyapunov spectra at each size:

\begin{center}
\begin{tabular}{lccc}
\toprule
 & $L=100$ & $L=200$ & $L=400$ \\
\midrule
True $\DKY$ & 22.14 & 45.46 & 89.66 \\
Model $\DKY$ (zero retraining) & 22.14 & 44.61 & 89.13 \\
Relative error & --- & 1.87\% & 0.60\% \\
\bottomrule
\end{tabular}
\end{center}

The error does not grow with extrapolation distance --- it shrinks. The
model's own $\DKY/L$ stays within 1\% of the true KS extensivity
constant ($\approx0.226$) at every size. To our knowledge this is the
first demonstration that a \emph{learned} latent representation
reproduces a chaotic PDE's own extensive scaling law at a domain size it
was never trained on, with literally zero parameter updates. Practically,
it means ensemble data assimilation at $L{=}400$ can reuse a filter built
once at $L{=}100$ --- an $\mathcal O(1)$ training cost instead of one
retraining per domain size. (\S\,3.2, \S\,5.1)

\section{Result 2: the mechanism does not survive the move to 2D}

The same architecture, retrained on each 2D system at a single domain
size, reproduces the true attractor well in-distribution (2D KS:
$\DKY{=}71.9$; Kolmogorov flow: $\DKY{=}21.3$ vs.\ true $23.3$). But
zero-shot transfer to a $1.5$--$2\times$ larger domain --- the exact
procedure that worked cleanly in 1D --- degrades sharply:

\begin{center}
\begin{tabular}{lcc}
\toprule
System & True $\DKY$ (transfer size) & Zero-shot error \\
\midrule
2D Kuramoto--Sivashinsky & 192.5 & \textbf{77\%} \\
Kolmogorov flow ($\alpha=0.05$ drag) & 56.8 & \textbf{50\%} \\
\bottomrule
\end{tabular}
\end{center}

Diagnosing \emph{why} separated two distinct failures. For 2D KS, the
true system's own dimension density $\DKY/\mathrm{Area}$ is confirmed
constant across two independently measured domain sizes ($0.0444$ vs.\
$0.0442$) --- the physics is genuinely extensive, so the 77\% gap is a
failure of the transfer \emph{mechanism} itself (our working hypothesis:
new domain area admits new pattern \emph{orientations} an isotropic
kernel's free weights were never shown during training, unlike 1D, where
new area only adds finer copies of an already-seen pattern). For
Kolmogorov flow, by contrast, the true $\DKY$ grew \emph{super-linearly}
with area ($2.44\times$ for a $2\times$ area increase) --- the physics
itself was not yet extensive at this drag strength, consistent with 2D
turbulence's well-known inverse energy cascade building a
domain-size-dependent large-scale structure when nothing caps it.
(\S\,8, \S\,8.1)

\section{Result 3: added linear drag restores extensivity --- and transfer}

That second diagnosis makes a specific, testable prediction: 2D
turbulence with linear (Ekman) drag has a known frictional arrest
scale, $\ell_\alpha\propto\alpha^{-3/2}$, above which the inverse
cascade is capped and the dynamics become extensive. If Kolmogorov
flow's transfer failure is a physics problem rather than an
architecture problem, \emph{simply increasing the drag, with no change
to the model at all, should repair it.} We doubled the drag
($\alpha:0.05\to0.10$), regenerated the dataset, and retrained the
identical architecture and recipe.

\begin{figure}[h]
\centering
\includegraphics[width=\textwidth]{figures/advisor_combined.pdf}
\caption{(a) Zero-shot domain-transfer error across every system/setting
tested --- doubling the drag alone takes Kolmogorov flow from 50\% down
to the same regime 1D KS achieves. (b) Why: true $\DKY/\mathrm{Area}$,
measured directly from the PDE (no learned model), is nearly flat out to
$2\times$ area at $\alpha=0.1$ (extensivity restored) before drifting
again at $3\times$ (delayed, not eliminated) --- and the model's own
transfer error tracks this drift almost exactly, rising back toward
$\sim$19\% exactly where the true density starts drifting again.}
\label{fig:combined}
\end{figure}

The result: transfer error at $2\times$ domain fell from 50\% to
\textbf{4.5\%}, confirmed independent of initial condition across three
seeds (4.5--5.9\%) --- ruling out a multistability artifact, the one
failure mode that produced a misleadingly large error estimate earlier
in this project for 2D KS. \textbf{A single physical parameter, chosen
for a reason derived from the physics, not tuned on the transfer metric
itself, recovered 1D-KS-quality transfer with zero architecture
changes.} One caveat, examined in the next section: of three
independently retrained propagators at $\alpha=0.1$, one collapsed under
transfer despite being indistinguishable from the other two at the
trained size. (\S\,8.1)

\section{Why some trained propagators transfer and others don't}
\label{sec:m2}

Three independently retrained propagators at $\alpha=0.1$ --- identical
architecture, dataset, and recipe, differing only in random seed --- give
a clean $2\times2$ puzzle. m0 and m1 transfer at $2\times$/$3\times$
domain size with 4.5--5.9\% error, matching \S4. \textbf{m2 collapses to
a near-fixed point under transfer ($D_{\mathrm{KY}}\approx0$--$2$,
independent of initial condition) while being statistically
indistinguishable from m0/m1 at the trained size} --- same training loss,
same native-size $D_{\mathrm{KY}}$. Neither the loss curve nor a
native-size evaluation would have flagged it.

\textbf{What differs, found without ever running a transfer test.} The
propagator's site-mixing layer has no softmax --- it is a literal,
linear, shift-invariant circular convolution over the lattice of latent
sites. For such an operator the gain at any spatial frequency is exact
and computable in closed form (the 2D Fourier transform of the kernel),
and training only ever constrains that gain at the handful of discrete
frequencies the \emph{trained} lattice size admits. Resizing to a larger
domain introduces new frequencies the kernel was never shown.

\begin{center}
\begin{tabular}{lccc}
\toprule
Model & Gain, trained freq.\ & Gain, new freq.\ ($2\times$) & Gain, new freq.\ ($3\times$) \\
\midrule
m0 & 2.02 & 1.58 & 1.82 \\
m1 & 2.12 & 1.49 & 1.78 \\
\textbf{m2} & \textbf{4.41} & \textbf{3.44} & \textbf{3.96} \\
\bottomrule
\end{tabular}
\end{center}

m2's kernel peaks at roughly \emph{double} the gain of m0/m1 at the
frequencies it was trained on, and --- unlike m0/m1, whose gain at
genuinely new frequencies drops noticeably below their own trained peak
--- m2's gain at those new frequencies stays close to its already-large
peak. Tellingly, the raw kernel size (Frobenius norm) is statistically
indistinguishable across all three models: m2's filter is not bigger, it
is more sharply tuned into one resonant mode, a difference invisible to
any norm-based or loss-based check. One honest gap: a direct 300-step
rollout of the base trajectory shows no visible difference between the
three models at any resize factor tested, so this gain analysis is a
plausible, measurement-consistent mechanism for the collapse, not yet a
proven one --- closing that gap (a short-horizon tangent-growth probe at
each resize factor) is the immediate next step. (\S\,9.1)

\textbf{What this might mean going forward.} If it holds up, it points
at a gap in how this whole class of model is trained and validated:

\begin{itemize}[nosep]
\item \emph{Training loss and native-size accuracy are not a
  certificate of out-of-distribution structural stability.} Two seeds
  that are indistinguishable by every metric we currently use to accept
  a checkpoint differ by $2\times$ in an operator property that, as far
  as we can tell, is the thing that actually matters for transfer.
\item \emph{The relevant property is cheaply, exactly computable for
  this whole architecture family} (any weight-shared, convolution-based
  local propagator), because the site-mixing layers are linear and
  shift-invariant by design. A maximum-gain bound, checked once after
  training and costing nothing at inference, is a plausible screen to
  catch a bad seed before ever running the expensive transfer-validation
  pipeline.
\item \emph{The contrast points at a specific regularizer, not just
  ``penalize the norm.''} m0/m1 and m2 have statistically identical bulk
  kernel size --- what differs is that m2's gain is sharply
  \emph{concentrated} into one resonant frequency rather than spread out.
  That is the actual failure mode a regularizer should target: penalize
  peakedness/concentration of the kernel's frequency response (e.g.\ its
  gain variance across frequency, or simply its max-to-mean gain ratio),
  not overall gain, which a blanket spectral-norm cap would suppress
  indiscriminately and likely cost real expressivity for no reason. A
  smoothly-varying response is also the natural thing to want
  independent of transfer: a resonant filter disproportionately amplifies
  whatever noise or error happens to land near its one sensitive
  frequency, so this plausibly improves in-distribution robustness too,
  not only transfer --- untested, and the natural next experiment.
\item \emph{This is a second, model-side precondition for transfer,
  alongside the physics-side one in \S4.} \S4 showed transfer needs the
  \emph{true system} to be extensive at the target size, checkable by
  measuring $D_{\mathrm{KY}}/\mathrm{Area}$ with no learned model
  involved. This section suggests transfer also needs the \emph{learned
  operator} to have bounded gain at frequencies it was never shown ---
  checkable, in principle, with no transfer test either. A general
  transfer-readiness checklist for this architecture family may end up
  being exactly these two checks, one on the physics and one on the
  trained weights.
\end{itemize}

\section{Implications and next steps}

\textbf{A general design principle.} Before investing in a local-field
model for an unfamiliar extensive-chaos system, measure true
$\DKY/\mathrm{Area}$ at two domain sizes --- cheap, no learned model
required. If it is flat, zero-shot transfer should work; if not, check
whether a physically motivated dissipation term can restore it before
assuming the architecture needs to change.

\textbf{A diagnostic for the onset of large-scale structure.} Because
the local surrogate can only represent dynamics built from a bounded
interaction range, its transfer error is a direct readout of how much of
a system's attractor dimension is \emph{non-local} at a given scale ---
and, in Fig.~\ref{fig:combined}b, that readout tracked the true extensivity drift almost
1:1 at both points measured. The same idea applied to a cloud-resolving
model in radiative-convective equilibrium would, in principle, locate
the domain-size threshold for convective self-aggregation --- a known,
actively studied transition in atmospheric science --- without an
aggregation index defined by hand, and would give a causal test
(disabling rain evaporation is known to remove the threshold's
domain-size dependence, exactly analogous to raising $\alpha$ here).
We checked the literature specifically for this combination (local
surrogate + zero-shot transfer error + Lyapunov validation, used as a
locality/extensivity diagnostic); we did not find it. (\S\,9)

\textbf{Concrete next steps, in order:}
\begin{enumerate}[nosep]
\item \emph{Error bars.} A multi-seed re-measurement is underway: the
  true $\DKY$ itself carries $\sim$9\% seed-to-seed scatter at the sizes
  tested, comparable to several effects above, and must be quantified
  before any of this is submission-ready.
\item \emph{Map the crossover.} Sweep drag and domain size together and
  test the predicted $L^*\propto\alpha^{-3/2}$ scaling directly against
  the measured extensivity crossover --- turns a demonstration into a
  quantitative, falsifiable physics claim.
\item \emph{A cheap proxy.} Replace the expensive Lyapunov comparison
  with forward-only statistics (energy spectrum, autocorrelation) so the
  diagnostic is affordable at cloud-resolving-model scale.
\item \emph{Close the gain-vs-collapse gap (\S\ref{sec:m2}).} A
  short-horizon tangent-growth probe at each resize factor, to confirm
  the spectral-gain mechanism directly; if confirmed, test whether a
  peakedness penalty on the mixing kernel's frequency response
  eliminates the seed-dependent collapse and improves in-distribution
  robustness, without costing accuracy.
\item \emph{Rayleigh--B\'enard, then a cloud-resolving model} in
  radiative-convective equilibrium, as the direct test of the
  atmospheric application above.
\end{enumerate}

\textbf{Publication framing.} Result 1 alone is a clean, submittable
finding. Results 2--3 together are stronger: not ``this architecture
transfers,'' but ``transfer is governed by a checkable physical
precondition, and when it fails, the reason is diagnosable and, in one
case, fixable with a single physically motivated parameter change.''
Item 5, if it works, is the paper's headline application.

\end{document}
